\begin{cabstract}
当前人工智能系统多依赖数据中的统计相关性，缺乏可解释性与鲁棒性。因果发现旨在从观测数据中揭示变量间内在作用机制，是推动人工智能从“关联驱动”向“因果驱动”演进的关键。然而，现有因果发现方法多聚焦于有向无环图（DAG）结构学习，忽视了对结构因果模型中因果函数的刻画。传统深度学习求解器基于通用逼近定理设计，其稠密连接与黑箱特性使得学到的因果机制缺乏可解释性。此外，现有方法应对真实场景数据，并且模型难以支持精确的干预决策与反事实推理。鉴于此，本文引入柯尔莫哥洛夫 - 阿诺德网络（KAN），利用其符号化潜力与函数分解特性，研究基于 KAN 的因果结构与函数联合发现方法。

首先，提出基于 KAN 的连续优化因果发现算法（DAG-KAN）。针对现有方法重结构轻机制的局限，该算法构建轻量级单层 KAN 架构，实现因果结构与函数映射的端到端联合学习。基于 Kolmogorov-Arnold 表示定理，利用可学习 B 样条基函数显式拟合变量间非线性因果机制。设计适配 KAN 结构的节点独立性矩阵，通过边稀疏性约束学习因果邻接矩阵，并引入 L1正则化策略作用于样条系数，实现因果边的自动剪枝。在连续优化框架下加入无环性约束，确保学习到有向无环图。实验表明，该方法在加性模型数据上的结构学习精度优于主流方法，在高维复杂场景下结构汉明距离平均降低 22.6\%，且函数拟合均方误差达到 $10^{-3}$ 精度。

其次，提出基于 KAN 的知识融合因果发现算法（Causal-KAN）。针对真实场景普遍存在的非参数或半参数数据，设计知识嵌入层与局部连接激活层架构。通过在输入层引入因果掩码机制，将部分已知因果父节点关系硬编码为网络约束，在保障因果忠实性的同时具有深层网络函数逼近能力。这种结构化先验有效缩小了模型搜索空间，避免非因果变量的干扰。实验表明，该方法通过嵌入因果知识可有效防止过拟合，且学习到的因果函数具有可解释性，对节点规模和噪声水平敏感度低，鲁棒性较好。

本文提出的 DAG-KAN 与 Causal-KAN 方法，将因果发现研究从因果存在性判断推进至因果作用机制解析，构建了兼具可解释性与实用鲁棒性的因果函数学习方法。该方法能够显式表达变量间的作用机制，为精准医疗、政策干预等需要理解干预效果的领域提供了可操作的因果知识。
\end{cabstract}
\ckeywords{因果发现；Kolmogorov-Arnold 网络；因果函数；知识融合}

\begin{eabstract}
Current artificial intelligence systems largely rely on statistical correlations in data, lacking interpretability and robustness. Causal discovery aims to uncover the intrinsic mechanisms underlying the relationships between variables in observational data, and is key to driving the evolution of artificial intelligence from a “correlation-driven” to a “causal-driven” paradigm. However, existing causal discovery methods primarily focus on learning directed acyclic graph (DAG) structures, neglecting the characterization of causal functions within structural causal models. Traditional deep learning solvers are designed based on general approximation theorems; their dense connections and black-box nature result in learned causal mechanisms that lack interpretability. Furthermore, existing methods struggle with real-world data, and the models often fail to support precise intervention decisions and counterfactual reasoning. In light of this, this paper introduces the Kolmogorov–Arnold Network (KAN) and leverages its symbolic potential and function decomposition capabilities to investigate a KAN-based method for the joint discovery of causal structures and functions.

First, we propose a continuous optimization-based causal discovery algorithm (DAG-KAN) based on KAN. To address the limitation of existing methods that prioritize structure over mechanism, this algorithm constructs a lightweight single-layer KAN architecture to achieve end-to-end joint learning of causal structure and function mapping. Based on the Kolmogorov-Arnold representation theorem, we use learnable B-spline basis functions to explicitly fit nonlinear causal mechanisms between variables. We design a node independence matrix tailored to the KAN structure, learn the causal adjacency matrix under edge sparsity constraints, and introduce an L1 regularization strategy applied to spline coefficients to achieve automatic pruning of causal edges. By incorporating acyclicity constraints within the continuous optimization framework, we ensure the learned graph is a directed acyclic graph (DAG). Experiments demonstrate that this method achieves higher structural learning accuracy than mainstream methods on additive model data, reduces the average structural Hamming distance by 22.6\% in high-dimensional complex scenarios, and attains a mean squared error of $10^{-3}$ in function fitting.

Second, we propose a knowledge-fusion causal discovery algorithm based on KAN (Causal-KAN). To address the prevalence of non-parametric or semi-parametric data in real-world scenarios, we design an architecture comprising a knowledge embedding layer and a locally connected activation layer. By introducing a causal masking mechanism at the input layer, we hard-code known causal parent-child relationships as network constraints, thereby ensuring causal fidelity while maintaining the function approximation capabilities of deep neural networks. This structured prior effectively reduces the model search space and avoids interference from non-causal variables. Experiments demonstrate that this method effectively prevents overfitting by embedding causal knowledge, and the learned causal functions are interpretable, exhibit low sensitivity to node scale and noise levels, and demonstrate good robustness.

The DAG-KAN and Causal-KAN methods proposed in this paper advance causal discovery research from determining the existence of causality to analyzing causal mechanisms, establishing a causal function learning method that combines interpretability with practical robustness. This method can explicitly express the mechanisms of action between variables, providing actionable causal knowledge for fields such as precision medicine and policy intervention that require an understanding of intervention effects.

\end{eabstract}
\ekeywords{Causal Discovery; Kolmogorov-Arnold Networks; Causal Function; Knowledge Fusion}

